% ============================================================
%  CAPÍTULO 3: MARCO TEÓRICO
% ============================================================
\chapter{Marco teórico}
\label{cap:teoria}

% ------------------------------------------------------------
\section{Introducción}
% ------------------------------------------------------------

El presente capítulo desarrolla el marco teórico que sustenta el análisis empírico
de la demanda de seguros sanitarios privados en España. Partimos de la premisa
fundamental de que los hogares españoles operan en un sistema sanitario dual,
donde la sanidad pública universal coexiste ---y compite--- con una oferta privada
creciente. La decisión de contratar un seguro privado no obedece únicamente a
consideraciones de renta, sino que incorpora una evaluación implícita de los
costes de oportunidad asociados al uso del sistema público, especialmente aquellos
derivados de los tiempos de espera.

Estructuramos el capítulo en tres bloques analíticos. En primer lugar,
formalizamos el modelo de decisión del hogar entre sanidad pública y privada,
identificando las variables clave que determinan la asignación del gasto sanitario.
En segundo lugar, introducimos la Ley de Little como fundamento teórico para la
elección de los días de espera como \textit{proxy} de saturación del sistema
público. Finalmente, desarrollamos el concepto de precio sombra del tiempo de
espera.

% ------------------------------------------------------------
\section{El modelo de decisión del hogar: sanidad pública versus privada}
% ------------------------------------------------------------

\subsection{Planteamiento general}

Consideremos un hogar representativo que maximiza su utilidad derivada del consumo
de bienes y servicios, incluyendo servicios sanitarios. El hogar dispone de acceso
gratuito en el punto de uso a la sanidad pública (financiada mediante impuestos),
pero puede optar por complementar o sustituir dicho acceso mediante la contratación
de un seguro sanitario privado.

Siguiendo la tradición de la economía de la salud iniciada por
\textcite{Grossman1972}, el individuo no demanda servicios sanitarios
\textit{per se}, sino \textit{salud} como bien de inversión y consumo. Los
servicios sanitarios son \textit{inputs} en la función de producción de salud del
hogar. En este contexto, la elección entre sanidad pública y privada puede
modelizarse como una decisión de asignación de recursos escasos ---tiempo y
dinero--- para maximizar el stock de salud esperado.

\subsection{Formalización del modelo}

Sea $U = U(C, H)$ la función de utilidad del hogar, donde $C$ representa el
consumo de bienes no sanitarios y $H$ el stock de salud. El hogar enfrenta la
restricción presupuestaria:

\begin{equation}
  Y = C + p_s \cdot S
  \label{eq:restriccion}
\end{equation}

donde $Y$ es la renta disponible, $p_s$ es la prima del seguro sanitario privado
y $S \in \{0,1\}$ es una variable indicadora de contratación del seguro.

La producción de salud depende de los servicios sanitarios consumidos:

\begin{equation}
  H = H(Q_{\text{pub}}, Q_{\text{priv}}, Z)
  \label{eq:produccion_salud}
\end{equation}

donde $Q_{\text{pub}}$ son los servicios obtenidos del sistema público,
$Q_{\text{priv}}$ los del sistema privado (accesibles solo si $S=1$), y $Z$ son
otros determinantes de la salud (estilo de vida, genética, etc.).

El elemento crucial de nuestro modelo es que el acceso a la sanidad pública,
aunque gratuito en términos monetarios, conlleva un \textbf{coste temporal} no
trivial:

\begin{equation}
  Q_{\text{pub}} = Q_{\text{pub}}(W)
  \label{eq:demanda_publica}
\end{equation}

donde $W$ representa el tiempo de espera medio para acceder a los servicios
públicos. A mayor $W$, menor es la cantidad efectiva de servicios que el hogar
puede consumir en un período dado.

\subsection{La decisión de contratación}

El hogar contratará el seguro privado ($S=1$) si y solo si:

\begin{equation}
  U\!\left(Y - p_s,\, H(Q_{\text{pub}}(W), Q_{\text{priv}}, Z)\right)
  > U\!\left(Y,\, H(Q_{\text{pub}}(W), 0, Z)\right)
  \label{eq:decision}
\end{equation}

Aplicando una aproximación de Taylor de primer orden, la condición puede
reescribirse como:

\begin{equation}
  \frac{\partial U}{\partial H} \cdot
  \bigl[ H(Q_{\text{pub}}, Q_{\text{priv}}, Z)
       - H(Q_{\text{pub}}, 0, Z) \bigr]
  > \frac{\partial U}{\partial C} \cdot p_s
  \label{eq:condicion_contratacion}
\end{equation}

Esta expresión revela que la decisión de contratar depende de: (i) la valoración
marginal de la salud respecto al consumo, (ii) el incremento de salud esperado por
acceder a la sanidad privada, y (iii) el precio del seguro. Crucialmente, cuando
$W$ aumenta, $Q_{\text{pub}}(W)$ disminuye, lo que incrementa el diferencial de
salud esperado y, por tanto, la probabilidad de contratar el seguro. Este es el
\textbf{mecanismo de huida} que fundamenta nuestra hipótesis empírica.

\subsection{Heterogeneidad según renta}

El modelo predice que el efecto de los tiempos de espera sobre la demanda de
seguros privados será heterogéneo según el nivel de renta del hogar. Para hogares
de renta baja, la restricción presupuestaria es más estricta y
$\partial U / \partial C$ es relativamente alto, lo que dificulta la contratación
incluso ante tiempos de espera elevados. Para hogares de renta alta, el coste de
oportunidad del tiempo es mayor, lo que amplifica el efecto de $W$ sobre la
decisión.

% ------------------------------------------------------------
\section{La Ley de Little y los tiempos de espera como indicador de saturación}
% ------------------------------------------------------------

\subsection{Fundamentos de teoría de colas}

La Ley de Little \parencite{Little1961} constituye uno de los resultados más
robustos y generales de la teoría de colas. Establece una relación matemática
fundamental entre tres magnitudes de cualquier sistema de espera en equilibrio
estacionario:

\begin{equation}
  L = \lambda \cdot W
  \label{eq:little}
\end{equation}

donde $L$ es el número medio de pacientes en el sistema, $\lambda$ es la tasa
media de llegadas (demanda de servicios por unidad de tiempo) y $W$ es el tiempo
medio de permanencia en el sistema. La elegancia de esta ley reside en su
generalidad: es válida independientemente de la distribución de llegadas, el
número de servidores o la disciplina de la cola.

\subsection{Aplicación al sistema sanitario público}

En el contexto del SNS español:

\begin{itemize}
  \item $L$: número de pacientes en lista de espera para una especialidad.
  \item $\lambda$: tasa de derivaciones desde atención primaria.
  \item $W$: días medios de espera para consulta con especialista.
\end{itemize}

\subsection{\texorpdfstring{¿Por qué \(W\) y no \(L\) como variable de saturación?}{Por que W y no L como variable de saturacion}}

Argumentamos que $W$ (tiempo de espera) es la variable relevante para el hogar,
no $L$ (número de pacientes en cola), por cuatro razones:

\begin{enumerate}
  \item El hogar no observa directamente $L$, pero sí experimenta $W$ cuando
    solicita una cita con el especialista.
  \item $W$ tiene una interpretación directa en términos de coste de oportunidad:
    90 días de espera implican tres meses de dolor, incertidumbre o incapacidad
    laboral.
  \item $W$ es comparable entre CCAA con independencia de su tamaño poblacional.
  \item Los individuos procesan la información en términos de tiempos más que de
    volúmenes \parencite{Kahneman1979}.
\end{enumerate}

% ------------------------------------------------------------
\section{El precio sombra del tiempo de espera}
% ------------------------------------------------------------

\subsection{Concepto de precio sombra}

El \textbf{precio sombra} de un recurso es el coste de oportunidad implícito de
su uso cuando dicho recurso no se intercambia en un mercado con precios explícitos.
En el caso de la sanidad pública, acceder a los servicios requiere invertir tiempo,
cuyo coste de oportunidad constituye el precio sombra del acceso.

\subsection{Formalización}

Sea $p_{\text{pub}}^{*}$ el precio sombra total de la sanidad pública para el
hogar:

\begin{equation}
  p_{\text{pub}}^{*} = c_{\text{directo}} + w \cdot T
  \label{eq:precio_sombra}
\end{equation}

donde $c_{\text{directo}}$ son los costes directos (copagos, transporte), $w$ es
el coste de oportunidad del tiempo (salario/hora) y $T$ es el tiempo total
invertido en acceder al servicio:

\begin{equation}
  T = W_{\text{cita}} + W_{\text{sala}} + T_{\text{desplazamiento}}
      + T_{\text{consulta}}
  \label{eq:tiempo_total}
\end{equation}

De estos componentes, $W_{\text{cita}}$ puede superar los 100 días en
especialidades saturadas como traumatología, oftalmología o dermatología.

\subsection{Efecto sobre el precio relativo}

El precio relativo de la sanidad privada respecto a la pública es:

\begin{equation}
  \pi = \frac{p_s + p_{\text{priv}}^{*}}{p_{\text{pub}}^{*}}
  \label{eq:precio_relativo}
\end{equation}

Cuando $W_{\text{cita}} \uparrow$, se tiene que $p_{\text{pub}}^{*} \uparrow$ y,
por tanto, $\pi \downarrow$: la sanidad privada se vuelve relativamente más barata.
Este abaratamiento relativo estimula la demanda de seguros privados.

\subsection{Elasticidad de sustitución}

La existencia de un precio sombra significativo implica que la sanidad pública y
privada son \textbf{bienes sustitutivos}. Definiendo la elasticidad de sustitución
como:

\begin{equation}
  \sigma = \frac{\partial \ln(Q_{\text{priv}}/Q_{\text{pub}})}
                {\partial \ln(\pi)}
  \label{eq:elasticidad_sustitucion}
\end{equation}

nuestra hipótesis empírica es que $\sigma < 0$: en términos operativos,
$\partial \ln(\text{gasto\_privado}) / \partial W > 0$.

% ------------------------------------------------------------
\section{Síntesis y proposiciones teóricas}
% ------------------------------------------------------------

Del marco teórico desarrollado se derivan las siguientes proposiciones que guían
el análisis empírico:

\begin{description}
  \item[\textbf{P1} (Efecto renta):] El gasto en seguros sanitarios privados es
    una función creciente de la renta disponible del hogar, con elasticidad-renta
    esperada mayor que la unidad.

  \item[\textbf{P2} (Efecto precio sombra):] El gasto en seguros privados es una
    función creciente de los tiempos de espera del sistema público.

  \item[\textbf{P3} (Heterogeneidad):] El efecto de los tiempos de espera sobre
    la demanda de seguros privados es mayor en las CCAA con mayor renta per cápita,
    debido al mayor coste de oportunidad del tiempo.

  \item[\textbf{P4} (Respuesta diferida):] El ajuste de la demanda de seguros
    privados a cambios en los tiempos de espera no es instantáneo; el efecto
    completo se materializa con un rezago de uno a dos años.
\end{description}

Estas proposiciones serán contrastadas empíricamente en el
Capítulo~\ref{cap:resultados} mediante modelos de panel que explotan la variación
temporal y transversal de las CCAA españolas en el período 2016--2024.
